% Reconstructed text from V5 PDF page 16; non-native source
\par\smallskip\noindent\rule{\textwidth}{0.4pt}\par{\sffamily\footnotesize Original V5 page 16 -- recovered text}\par\begin{flushleft}\footnotesize\ttfamily
\par\smallskip
\par\smallskip
\noindent Diagnostics and Serving Infrastructure for MAS 16\par
\par\smallskip
\par\smallskip
\par\smallskip
\noindent mal causal estimators; the method requires open-weight SLMs (closed APIs are incompatible\par
\noindent with prefill-stage signal extraction); only failed traces are evaluated, with false-positive control\par
\noindent on mixed-success runs left as future work; the symptom ratio must be manually configured per\par
\noindent benchmark; and on short traces (the AG subset of Who\&When), explicit reasoning methods\par
\noindent such as A2P outperform MASPrism (43.65\% vs. 36.51\%, Table 1). Without the filtering stage,\par
\noindent the method encounters out-of-memory failures on 73 of 117 GAIA traces (Table 3).\par
\par\smallskip
\noindent Additional LLM-guided methods. Beyond the four systems detailed above, the LLM-guided\par
\noindent family has expanded rapidly. RAFFLES [51] applies iterative reasoning chains to progressively\par
\noindent narrow the failure location, achieving 50.79\% step-level accuracy on algorithm-generated sce-\par
\noindent narios in the Who\&When benchmark (EACL 2026). VerifyMAS [31] formulates attribution as\par
\noindent hypothesis verification using natural language inference (NLI), optionally fine-tuned via SFT,\par
\noindent producing agent-level attribution paired with error-type classification.\par
\noindent  At the abstract-verified tier, additional systems reveal design directions: CausalFlow applies\par
\noindent causal responsibility scores to agent trajectories via step-level counterfactual intervention. CAR\par
\noindent (Causal Agent Replay) combines structural causal models with five do-operations and Shapley\par
\noindent credit-splitting for multi-agent attribution. DCFA (EMNLP 2026) introduces a training-free\par
\noindent dual-view approach: a global causal dependency graph combined with local counterfactual rea-\par
\noindent soning, achieving +8.27\% step accuracy over the state of the art. HPFA (COLM 2026) uses\par
\noindent hypergraph-based paired failure attribution with SFT + RL training. EDGE (EMNLP 2026)\par
\noindent constructs multi-error dependency graphs, addressing the challenging case of multiple simul-\par
\noindent taneous failures. DeFA achieves the highest accuracy on both Who\&When and Who\&When\par
\noindent Pro’s text subset through protocol + semantic dependency propagation. EMFA establishes a\par
\noindent new state of the art on step-level Who\&When attribution by modeling cascading + persistent\par
\noindent error propagation loops.\par
\noindent  These systems illustrate the field’s rapid progression from single-failure linear analysis to-\par
\noindent ward multi-fault, graph-structured attribution. Three additional architectural directions are visi-\par
\noindent ble: SDBL (AAAI 2026) proposes scoping the failure before localization — defining the failure\par
\noindent boundary first, then identifying the root cause within that boundary. DuoTrace separates detec-\par
\noindent tion from attribution, using a first pass to identify whether a failure occurred and a second pass\par
\noindent to attribute where it originated. REFLECT (ICML 2026 Workshop) bridges diagnosis to inter-\par
\noindent vention by generating repair actions from attributed failures — the only system in our survey\par
\noindent that closes the loop from attribution to recovery in a single framework. Together, these 23 addi-\par
\noindent tional systems (all abstract-verified) suggest that the failure-localization landscape is maturing\par
\noindent from isolated tools toward integrated pipelines with standardized interfaces.\par
\par\smallskip
\noindent Benchmarks and evaluation. TRAIL [10] provides the first OTel-format benchmark for trace\par
\noindent debugging: 1,987 spans across 148 traces from GAIA (multi-agent) and SWE-Bench Lite\par
\noindent (single-agent) systems, with structured error categorization. The best model (Gemini 2.5 Pro)\par
\noindent achieves only ∼11\% joint accuracy, revealing that current LLMs struggle with long structured\par
\noindent traces. Token consumption is substantial: GAIA traces average 263–314K input tokens, while\par
\noindent SWE-Bench traces average 617–728K, with a pronounced anti-correlation between trace length\par
\noindent and accuracy (Pearson r = −0.29 to −0.38, §5.1.3). The annotation protocol required \$12.66\par
\noindent per trace and 30–40 minutes of expert time, illustrating the cost of ground-truth creation for\par
\noindent diagnostic evaluation.\par
\noindent  RAFFLES [51] achieves 50.79\% step-level accuracy on algorithm-generated scenarios through\par
\noindent iterative reasoning (K=2 rounds) at \$0.12 per short log and \$0.31 per long log using Claude\par
\noindent Sonnet — 21–41× cheaper than a junior engineer (§5.1, Appendix D). Parallelized latency is\par
\end{flushleft}
